\begin{theorem}[The core simulation]
  \label{thm:aba-coresim}
  \lean{PLTS.ABA.coreSim}\leanok
  \uses{def:aba-hybrids, def:aba-spec}
  The analysis-side composition is forward-simulated by the repaired ABA specification under the relation $Inv \wedge Abs$.
  The abstraction $Abs$ pins the corrupted set $F$, the returns and $coin = \bot$, so the twin never fires the coin rule, and it is two-phase on the twin's $val$ (D16): before the first return the twin is unbound, its abstract calls equalling the concrete write-once inputs and its ghost synced; after it the board is clear and $bind = val = some\,v$ with $v$ certified by a concrete A-lock.
  Every hidden row and every concrete coin flip is answered by a stutter under a constant coupling, the only burst being $\mathrm{decide\_burst}$ at the first $\mathrm{retABA}$, which within one weak step binds, fills, decides and returns.
  The concrete invariant carries the cross-round Graded Agreement machinery, downward C-lock propagation and upward bind commitment.
\end{theorem}

\begin{proof}\leanok
  \uses{def:aba-hybrids, def:aba-spec, def:weakstep}
  The relation is checked row by row, one step class at a time, with $Inv$ preserved at each.
  Every hidden row and the coin flip are answered by an $Abs$-frame stutter, the twin holding still under a constant coupling; since $Abs$ pins $coin = \bot$ the twin never fires the coin rule, so the probabilistic row couples to that stutter.
  The first $\mathrm{retABA}$ is answered by the single $\mathrm{decide\_burst}$ weak step, whose bind is licensed by the $f{+}1$ input support carried in the invariant and whose $A$-lock certificate comes from a pigeonhole on the $n-f$ DECIDED tally behind the return.
  A $\mathrm{fail}$ row is answered by the corruption broadcast.
\end{proof}
